\documentclass[11pt,a4paper]{article}

% ── Codificación, fuentes y lenguaje ───────────────────────────────────────────
\usepackage[utf8]{inputenc}
\usepackage[T1]{fontenc}
\usepackage[default,scale=0.95]{sourcesanspro}
\renewcommand{\familydefault}{\sfdefault}
\usepackage[spanish,es-noshorthands,es-tabla]{babel}

% ── Geometría y espaciado ─────────────────────────────────────────────────────
\usepackage[top=2.2cm, bottom=2.5cm, left=2.2cm, right=2.2cm, headheight=15pt]{geometry}
\usepackage{setspace}
\setstretch{1.18}
\usepackage{parskip}

% ── Matemáticas y unidades ────────────────────────────────────────────────────
\usepackage{amsmath}
\usepackage{amssymb}
\usepackage{siunitx}

% ── Circuitos ─────────────────────────────────────────────────────────────────
\usepackage[european, straightvoltages]{circuitikz}

% ── Tablas ────────────────────────────────────────────────────────────────────
\usepackage{booktabs}
\usepackage{tabularx}
\usepackage{array}
\usepackage{longtable}
\usepackage{colortbl}
\usepackage{enumitem}

% ── Colores, cajas e iconos ───────────────────────────────────────────────────
\usepackage[dvipsnames]{xcolor}
\usepackage{tcolorbox}
\tcbuselibrary{skins, breakable}
\usepackage{fontawesome5}
\usepackage{graphicx}

% ── Listados de código (dark-mode infográfico) ────────────────────────────────
%   Estilo base: fondo oscuro con números de línea. Los bloques lstlisting se
%   envuelven automáticamente en una caja tcolorbox redondeada.
\usepackage{listings}

% ── Secciones con barra de color (infografía) ────────────────────────────────
\usepackage{titlesec}
\titleformat{\section}
  {\normalfont\LARGE\bfseries\color{azulNoche}}
  {\colorbox{cyanNeon}{\color{fondoOscuro}\thesection}}{0.7em}{}
  [\vspace{2pt}\color{cyanNeon}\rule{\linewidth}{1.8pt}]
\titlespacing*{\section}{0pt}{24pt}{10pt}
\titleformat{\subsection}
  {\normalfont\large\bfseries\color{azulNoche}}
  {\textcolor{cyanNeon}{\thesubsection}}{0.7em}{}
  [\vspace{1pt}\color{grisLinea}\rule{\linewidth}{0.6pt}]
\titlespacing*{\subsection}{0pt}{14pt}{6pt}
\titleformat{\subsubsection}
  {\normalfont\normalsize\bfseries\color{azulNoche}}
  {\textcolor{cyanNeon}{\thesubsubsection}}{0.7em}{}
\titlespacing*{\subsubsection}{0pt}{10pt}{4pt}

% ── Cabeceras y pies de página ────────────────────────────────────────────────
\usepackage{fancyhdr}
\usepackage{lastpage}
\pagestyle{fancy}
\fancyhf{}
\renewcommand{\headrulewidth}{0pt}
\renewcommand{\footrulewidth}{0pt}
\fancyfoot[C]{%
  \begin{minipage}{\linewidth}
    \vspace{2pt}
    \color{cyanNeon}\rule{\linewidth}{1.2pt}\\[2pt]
    \centering\color{grisTexto}\footnotesize
    Página \thepage\ de \pageref{LastPage}
  \end{minipage}%
}

% ── Hiperenlaces ──────────────────────────────────────────────────────────────
\usepackage[colorlinks=true, linkcolor=azulNoche, urlcolor=cyanNeon]{hyperref}
\sloppy
\emergencystretch 3em

% ── Paleta de colores — tecnológico dark-mode ──────────────────────────────────
\definecolor{fondoOscuro}{HTML}{0D1B2A}%             Dark navy — fondo banda título
\definecolor{verdeTurquesa}{HTML}{004D40}%           Verde turquesa — bandas de marca
\definecolor{azulNoche}{HTML}{1B3A6B}%               Midnight blue — secciones, marcos
\definecolor{azulMedio}{HTML}{2A5298}%               Azul medio — variante de acento
\definecolor{cyanNeon}{HTML}{00C8E0}%                Cyan eléctrico — acentos primarios
\definecolor{verdeSignal}{HTML}{00B887}%             Verde señal — fortalezas, éxito
\definecolor{naranjaVivo}{HTML}{FF7043}%             Naranja vivo — mejoras, advertencias
\definecolor{rojoAlerta}{HTML}{E53935}%              Rojo alerta — errores
\definecolor{grisPapel}{HTML}{F4F6FA}%               Fondo tarjetas claras
\definecolor{grisLinea}{HTML}{C8D0E0}%               Bordes sutiles
\definecolor{grisTexto}{HTML}{6B7A99}%               Texto secundario
\definecolor{amarilloNota}{HTML}{FFB300}%            Notas / advertencias doradas
\definecolor{codigoFondo}{HTML}{1E2D3D}%             Fondo del bloque de código
\definecolor{codigoComentario}{HTML}{7A8FA6}%        Comentarios
\definecolor{codigoCadena}{HTML}{FF9D5C}%            Cadenas / strings
\definecolor{codigoKeyword}{HTML}{00C8E0}%           Palabras clave
\definecolor{codigoNumero}{HTML}{00E5A0}%            Números

% ── Banda de título: portada infográfica ──────────────────────────────────────
%   Diseño en 2 capas: banda de color (por defecto fondo oscuro) + franja de
%   acento cyan abajo. Uso: \bandaTitulo[color]{título}{subtítulo}.
%   El icono se puede cambiar con \renewcommand{\iconoBanda}{...} (FontAwesome).
\newcommand{\iconoBanda}{microchip}
\newcommand{\bandaTitulo}[3][fondoOscuro]{%
  \begin{tcolorbox}[
    enhanced,
    colback=#1, colframe=#1,
    boxrule=0pt, arc=8pt, outer arc=8pt,
    width=\linewidth,
    top=18pt, bottom=6pt, left=14pt, right=14pt,
    borderline south={3.5pt}{0pt}{cyanNeon},
  ]
    \begin{minipage}[c]{0.07\linewidth}
      \centering
      \textcolor{cyanNeon}{\fontsize{30}{30}\selectfont\faIcon{\iconoBanda}}
    \end{minipage}%
    \hspace{8pt}%
    \begin{minipage}[c]{0.88\linewidth}
      {\fontsize{20}{24}\selectfont\bfseries\color{white}#2}\par\vspace{5pt}%
      \textcolor{cyanNeon!80!white}{\small\faIcon{angle-right}~#3}%
    \end{minipage}
    \vspace{4pt}
  \end{tcolorbox}%
  \vspace{8pt}%
}

% ── Tarjetas de datos (infografía) ────────────────────────────────────────────
\tcbset{
  tarjetaDato/.style={
    enhanced, breakable,
    arc=6pt, outer arc=6pt,
    colback=grisPapel, colframe=azulNoche,
    boxrule=1pt,
    fonttitle=\bfseries\small, coltitle=white,
    attach boxed title to top left={yshift=-3mm, xshift=6mm},
    boxed title style={
      colback=azulNoche, colframe=azulNoche,
      arc=4pt, boxrule=0pt,
    },
    drop fuzzy shadow=azulNoche!30!white,
    top=8pt, bottom=8pt, left=10pt, right=10pt,
  },
  tarjetaIcono/.style={
    enhanced, breakable,
    arc=6pt, outer arc=6pt,
    colback=white, colframe=grisLinea, boxrule=0.8pt,
    fonttitle=\bfseries\small, coltitle=azulNoche,
    borderline west={3pt}{0pt}{cyanNeon},
    drop fuzzy shadow=azulNoche!25!white,
    top=6pt, bottom=6pt, left=10pt, right=10pt,
  },
  cajaTitulo/.style={
    enhanced, breakable,
    arc=6pt, outer arc=6pt,
    colback=azulNoche!10!grisPapel, colframe=azulNoche,
    boxrule=1pt,
    fonttitle=\bfseries, coltitle=white,
    attach boxed title to top left={yshift=-3mm, xshift=6mm},
    boxed title style={
      colback=azulNoche, colframe=azulNoche,
      arc=4pt, boxrule=0pt,
    },
    drop fuzzy shadow=azulNoche!25!white,
    top=8pt, bottom=8pt, left=10pt, right=10pt,
  },
  cajaContenido/.style={
    enhanced, breakable,
    arc=5pt, outer arc=5pt,
    colback=grisPapel, colframe=grisLinea,
    boxrule=0.8pt,
    fonttitle=\bfseries\small, coltitle=azulNoche,
    top=6pt, bottom=6pt, left=10pt, right=10pt,
  },
  cajaFortaleza/.style={
    enhanced, breakable,
    arc=6pt, outer arc=6pt,
    colback=verdeSignal!8!white, colframe=verdeSignal,
    boxrule=1pt,
    borderline west={4pt}{0pt}{verdeSignal},
    fonttitle=\bfseries\small, coltitle=verdeSignal!70!black,
    drop fuzzy shadow=verdeSignal!20!white,
    top=6pt, bottom=6pt, left=10pt, right=10pt,
  },
  cajaMejora/.style={
    enhanced, breakable,
    arc=6pt, outer arc=6pt,
    colback=naranjaVivo!8!white, colframe=naranjaVivo,
    boxrule=1pt,
    borderline west={4pt}{0pt}{naranjaVivo},
    fonttitle=\bfseries\small, coltitle=naranjaVivo!80!black,
    drop fuzzy shadow=naranjaVivo!20!white,
    top=6pt, bottom=6pt, left=10pt, right=10pt,
  },
  cajaRecomendacion/.style={
    enhanced, breakable,
    arc=6pt, outer arc=6pt,
    colback=cyanNeon!6!white, colframe=cyanNeon!60!azulNoche,
    boxrule=1pt,
    borderline west={4pt}{0pt}{cyanNeon},
    fonttitle=\bfseries\small, coltitle=azulNoche,
    drop fuzzy shadow=azulNoche!20!white,
    top=6pt, bottom=6pt, left=10pt, right=10pt,
  },
}

% ── Ícono + texto (encabezados de sección y elementos) ────────────────────────
\newcommand{\iconotexto}[2]{\textcolor{cyanNeon}{\faIcon{#1}}~#2}

% ── Listados de código: estilo dark + auto-envoltura ─────────────────────────
\lstdefinestyle{estiloCodigo}{
    backgroundcolor=\color{codigoFondo},
    basicstyle=\ttfamily\small\color{white},
    commentstyle=\color{codigoComentario}\itshape,
    stringstyle=\color{codigoCadena},
    keywordstyle=\color{codigoKeyword}\bfseries,
    numberstyle=\tiny\color{codigoComentario},
    numbers=left,
    numbersep=10pt,
    showstringspaces=false,
    breaklines=true,
    breakatwhitespace=false,
    tabsize=4,
    frame=none,
    captionpos=b,
    aboveskip=6pt,
    belowskip=4pt,
    xleftmargin=14pt,
    framexleftmargin=14pt,
    literate=
      {á}{{\'a}}1 {é}{{\'e}}1 {í}{{\'i}}1 {ó}{{\'o}}1 {ú}{{\'u}}1
      {Á}{{\'A}}1 {É}{{\'E}}1 {Í}{{\'I}}1 {Ó}{{\'O}}1 {Ú}{{\'U}}1
      {ñ}{{\~n}}1 {Ñ}{{\~N}}1 {ü}{{\"u}}1 {Ü}{{\"U}}1
      {¿}{{?`}}1 {¡}{{!`}}1
}
\lstdefinestyle{estiloPython}{
    style=estiloCodigo,
    language=Python,
}
\lstset{style=estiloCodigo}
\BeforeBeginEnvironment{lstlisting}{%
  \begin{tcolorbox}[
    enhanced, breakable,
    arc=6pt, outer arc=6pt,
    colback=codigoFondo, colframe=azulNoche,
    boxrule=1pt,
    drop fuzzy shadow=azulNoche!25!white,
    top=2pt, bottom=2pt, left=0pt, right=0pt,
  ]%
}
\AfterEndEnvironment{lstlisting}{\end{tcolorbox}}

% ── Cajas de contenido temáticas ──────────────────────────────────────────────
\newtcolorbox{cajaRecuerda}{
    enhanced, breakable,
    arc=6pt, outer arc=6pt,
    colback=azulNoche!10!grisPapel,
    colframe=azulNoche, boxrule=1pt,
    borderline west={4pt}{0pt}{cyanNeon},
    fonttitle=\bfseries\small\color{white},
    title={\faIcon{info-circle}~Recuerda},
    attach boxed title to top left={yshift=-3mm, xshift=6mm},
    boxed title style={colback=azulNoche, arc=4pt, boxrule=0pt},
    drop fuzzy shadow=azulNoche!25!white,
    left=10pt, right=10pt, top=8pt, bottom=8pt,
}

\newtcolorbox{cajaConcepto}{
    enhanced, breakable,
    arc=6pt, outer arc=6pt,
    colback=verdeSignal!8!white,
    colframe=verdeSignal, boxrule=1pt,
    borderline west={4pt}{0pt}{verdeSignal},
    fonttitle=\bfseries\small\color{white},
    title={\faIcon{lightbulb}~Concepto clave},
    attach boxed title to top left={yshift=-3mm, xshift=6mm},
    boxed title style={colback=verdeSignal!80!black, arc=4pt, boxrule=0pt},
    drop fuzzy shadow=verdeSignal!20!white,
    left=10pt, right=10pt, top=8pt, bottom=8pt,
}

\newtcolorbox{cajaEjemplo}{
    enhanced, breakable,
    arc=6pt, outer arc=6pt,
    colback=cyanNeon!5!white,
    colframe=cyanNeon!70!azulNoche, boxrule=1pt,
    borderline west={4pt}{0pt}{cyanNeon},
    fonttitle=\bfseries\small\color{fondoOscuro},
    title={\faIcon{code}~Ejemplo},
    attach boxed title to top left={yshift=-3mm, xshift=6mm},
    boxed title style={colback=cyanNeon, arc=4pt, boxrule=0pt},
    drop fuzzy shadow=azulNoche!20!white,
    left=10pt, right=10pt, top=8pt, bottom=8pt,
}

\newtcolorbox{cajaEjercicio}{
    enhanced, breakable,
    arc=6pt, outer arc=6pt,
    colback=azulMedio!7!white,
    colframe=azulMedio, boxrule=1pt,
    borderline west={4pt}{0pt}{azulMedio},
    fonttitle=\bfseries\small\color{white},
    title={\faIcon{pencil-alt}~Ejercicio},
    attach boxed title to top left={yshift=-3mm, xshift=6mm},
    boxed title style={colback=azulMedio, arc=4pt, boxrule=0pt},
    drop fuzzy shadow=azulNoche!20!white,
    left=10pt, right=10pt, top=8pt, bottom=8pt,
}

\newtcolorbox{cajaReto}{
    enhanced, breakable,
    arc=6pt, outer arc=6pt,
    colback=naranjaVivo!6!white,
    colframe=naranjaVivo, boxrule=1pt,
    borderline west={4pt}{0pt}{naranjaVivo},
    fonttitle=\bfseries\small\color{white},
    title={\faIcon{fire}~Reto},
    attach boxed title to top left={yshift=-3mm, xshift=6mm},
    boxed title style={colback=naranjaVivo, arc=4pt, boxrule=0pt},
    drop fuzzy shadow=naranjaVivo!20!white,
    left=10pt, right=10pt, top=8pt, bottom=8pt,
}

% ── Indicadores de dificultad ─────────────────────────────────────────────────
\newcommand{\facil}{\textcolor{verdeSignal}{\faIcon{circle}\,\textbf{Fácil}}}
\newcommand{\intermedio}{\textcolor{naranjaVivo}{\faIcon{circle}\,\textbf{Intermedio}}}
\newcommand{\dificil}{\textcolor{rojoAlerta}{\faIcon{circle}\,\textbf{Difícil}}}

% ─────────────────────────────────────────────────────────────────────────────
\renewcommand{\iconoBanda}{atom}
\fancyhead[L]{\footnotesize\color{grisTexto}\faIcon{atom}~LABORATORIO\_II\_FISICA}
\fancyhead[R]{\footnotesize\color{grisTexto}Práctica 08}
\hypersetup{pdftitle={Laboratorio II de Física},pdfauthor={Prof. César Rodríguez}
  pdfsubject={LABORATORIO\_II\_FISICA}}

\begin{document}
\bandaTitulo[fondoOscuro]{Práctica 08 --- Carga y descarga de un capacitor con osciloscopio}{Semana 8 --- Fase 2: Elementos no lineales y régimen senoidal}

\begin{center}\small\color{grisTexto}Prof. César Rodríguez \textbar{} Universidad de Oriente, Núcleo de Sucre, Departamento de Física \textbar{} 2 de noviembre de 2026\end{center}

\begin{tabularx}{\linewidth}{@{}>{\raggedright\arraybackslash}p{4.6cm}X@{}}
\toprule
\textbf{Duración} & 3 horas de laboratorio presencial más 2 horas de análisis \\
\textbf{Modalidad} & 3 grupos con roles rotativos \\
\textbf{Entrega} & GitHub del grupo, carpeta semana\_08, antes del corte semanal \\
\textbf{Instrumento} & Práctica formativa con informe y control de fase \\
\bottomrule
\end{tabularx}

\section{\iconotexto{bullseye}{Objetivo}}

Familiarizarse con el uso del osciloscopio, visualizar las curvas de carga y descarga de un capacitor y determinar su capacitancia a partir de las medidas tomadas sobre esas curvas.

\section{\iconotexto{clipboard-list}{Contenidos que se aplican}}

\begin{itemize}[leftmargin=1.4em,itemsep=2pt,topsep=3pt]
  \item Régimen transitorio y estado estacionario.
  \item Ecuación diferencial del circuito RC y su solución exponencial.
  \item Constante de tiempo y semiperiodo.
  \item Linealización logarítmica de una exponencial.
  \item Energía almacenada en un capacitor y disipada en el resistor.
  \item Uso del osciloscopio en disparo único para capturar el transitorio.
  \item Resistencia serie equivalente incluyendo la del generador.
\end{itemize}

\section{\iconotexto{book-open}{Fundamento teórico}}

Aplicando la ley de tensiones al lazo de carga se obtiene la ecuación
diferencial cuya solución es exponencial:
\[
  v_C(t) = \varepsilon\left(1 - e^{-t/\tau}\right)
  \qquad\text{con}\qquad
  \tau = R C
\]
Durante la descarga, partiendo de $V_0$:
\[
  v_C(t) = V_0\,e^{-t/\tau}
\]
En un tiempo igual a $\tau$ el capacitor alcanza el 63,2 por ciento de la
tensión final, y el semiperiodo de la descarga vale
\[
  t_{1/2} = \tau \ln 2 \approx 0{,}693\,\tau
\]
Linealizando la descarga se obtiene una recta cuya pendiente es $-1/\tau$:
\[
  \ln v_C(t) = \ln V_0 - \frac{t}{\tau}
\]
Si los datos se toman en la carga conviene linealizar
$\ln\left[(V_0 - v_C)/V_0\right]$, que da la misma pendiente $-1/\tau$.
La energía almacenada al final de la carga es $W_C = \tfrac12 C \varepsilon^2$,
mientras que la fuente entrega $C\varepsilon^2$; la mitad se disipa en el
resistor independientemente del valor de $R$.
La resistencia que gobierna el transitorio no es solo la de la caja de
décadas: incluye la resistencia de salida del generador, de $50\ \Omega$ o
$600\ \Omega$ según el equipo, de modo que $R = R_{\text{década}} + R_{\text{generador}}$.

\begin{center}
\begin{circuitikz}
  \draw (0,0) to[V, l=$\varepsilon$] (0,3) to[closing switch, l=$S$] (2.4,3)
        to[R, l=$R$] (4.8,3) to[C, l=$C$] (4.8,0) -- (0,0);
  \draw (4.8,3) -- (6.4,3) to[voltmeter, l=$v_C$] (6.4,0) -- (4.8,0);
\end{circuitikz}
\par\smallskip{\small\itshape Circuito de carga y descarga; el conmutador permite alternar entre ambos transitorios}
\end{center}

\section{\iconotexto{box-open}{Materiales y equipos}}

\begin{itemize}[leftmargin=1.4em,itemsep=2pt,topsep=3pt]
  \item Un osciloscopio
  \item Un generador de funciones con onda cuadrada
  \item Un resistor de décadas
  \item Un capacitor del inventario del laboratorio
  \item Un óhmetro para verificar la resistencia conectada
\end{itemize}

\section{\iconotexto{flask}{Escenario}}

Un generador de onda cuadrada sustituye a la fuente y al interruptor, de modo que el capacitor se carga durante el tiempo de encendido y se descarga durante el de apagado. El grupo captura la curva en el osciloscopio y determina la capacitancia a partir de la pendiente de la recta que resulta de linealizar la exponencial.

\section{\iconotexto{list-ol}{Procedimiento}}

\begin{enumerate}[leftmargin=1.4em,itemsep=2pt,topsep=3pt]
  \item Armar el circuito RC con el generador de onda cuadrada y hacerlo revisar por el profesor.
  \item Ajustar las escalas de tiempo y de tensión del osciloscopio hasta tener en pantalla una sola curva de carga y descarga bien estable.
  \item Tomar un número suficiente de valores de $v_C$ para distintos tiempos, con la imagen estable.
  \item Medir con el óhmetro la resistencia conectada en serie con el capacitor y anotar la resistencia de entrada del generador.
  \item Medir $V_0$ conectando el generador en paralelo con el osciloscopio.
  \item Cambiar el capacitor por otro de valor unas diez veces mayor o menor, sin cambiar la frecuencia de la señal.
  \item Describir lo que se observa con el nuevo capacitor y explicar por qué cambia la curva.
\end{enumerate}

\section{\iconotexto{table}{Tabla de datos}}

\begin{center}
\begin{minipage}{\linewidth}
\centering
{\small
\textbf{Transitorio de carga y de descarga}\par\smallskip
\begin{tabular}{@{}cccc@{}}
\toprule
$t$ (s) & $v_C$ carga (V) & $v_C$ descarga (V) & $\ln\left[(V_0-v_C)/V_0\right]$ \\
\midrule
0,0 & 0,00 & $\ldots$ & $\ldots$ \\
0,5 & $\ldots$ & $\ldots$ & $\ldots$ \\
1,0 & $\ldots$ & $\ldots$ & $\ldots$ \\
1,5 & $\ldots$ & $\ldots$ & $\ldots$ \\
2,0 & $\ldots$ & $\ldots$ & $\ldots$ \\
3,0 & $\ldots$ & $\ldots$ & $\ldots$ \\
5,0 & $\ldots$ & $\ldots$ & $\ldots$ \\
\bottomrule
\end{tabular}
}\par\smallskip{\footnotesize El tiempo se toma desde el inicio del flanco; anotar $V_0$ y la resistencia total usada en el ajuste.}
\end{minipage}
\end{center}

\section{\iconotexto{calculator}{Análisis y tratamiento de datos}}

\begin{enumerate}[leftmargin=1.4em,itemsep=2pt,topsep=3pt]
  \item Graficar $\ln\left[(V_0-v_C)/V_0\right]$ frente a $t$ para la carga, y $\ln(v_C/V_0)$ si los datos son de la descarga.
  \item Ajustar la recta por mínimos cuadrados y determinar la pendiente con su error.
  \item Obtener la capacitancia a partir de la pendiente, tomando como $R$ la suma de la resistencia de décadas y la del generador (50 o 600 $\Omega$).
  \item Calcular la desviación porcentual del valor de C obtenido frente al valor del fabricante del apéndice de capacitores.
  \item Superponer la curva medida con el modelo exponencial y discutir los desvíos en el inicio del transitorio.
\end{enumerate}

\section{\iconotexto{plus-circle}{Ampliación (opcional)}}

\begin{itemize}[leftmargin=1.4em,itemsep=2pt,topsep=3pt]
  \item Calcular la energía almacenada y la entregada por la fuente y verificar el balance del 50 por ciento.
  \item Repetir el ajuste por mínimos cuadrados no lineales y comparar la constante de tiempo con la obtenida de la recta linealizada.
\end{itemize}

\section{\iconotexto{box-open}{Entregable}}

Curvas de carga y descarga capturadas, recta linealizada con su ajuste, capacitancia con su incertidumbre y desviación porcentual frente al valor del fabricante.

\section{\iconotexto{ruler-combined}{Rúbrica de corrección}}

\begin{tabularx}{\linewidth}{@{}X r@{}}
\toprule
Criterio & Peso\\
\midrule
Fundamento teórico y modelo aplicado & 1,0 \\
Montaje, seguridad y manejo del instrumental & 1,5 \\
Toma de datos completa y trazable & 2,0 \\
Presupuesto de incertidumbres & 2,0 \\
Análisis, ajuste y gráficas & 2,0 \\
Conclusiones contrastadas con el modelo & 1,0 \\
Presentación y estructura del informe & 0,5 \\
\midrule
\textbf{Total} & \textbf{10,0}\\
\bottomrule
\end{tabularx}

\section{\iconotexto{code}{Contrato del script}}

\begin{cajaConcepto}Todo script entregado imprime \textbf{un solo objeto JSON} por la salida estándar y devuelve un diccionario con estas claves. \texttt{status} vale \texttt{ok} o \texttt{error}; si es \texttt{error}, \texttt{errores} explica el motivo y el programa sale con código distinto de cero.
\begin{lstlisting}[language=Python,style=estiloPython]
{
  "status": "ok",
  "practica": 1,
  "magnitud": "resistencia",
  "valor": 99.4,
  "unidad": "ohm",
  "incertidumbre": 0.6,
  "metodo": "ajuste_lineal_minimos_cuadrados",
  "n_puntos": 8,
  "residual_max": 0.012,
  "fuente_datos": "semana_01/datos/iv_resistor_220.csv",
  "errores": []
}

\end{lstlisting}
\end{cajaConcepto}

\begin{tcolorbox}[cajaMejora]\textbf{Antes de empezar:} Antes de energizar, contrasta el montaje con la figura y verifica rangos y polaridad de cada instrumento. Si un instrumento se calienta, huele o marca fuera de escala, corta la alimentación y avisa al docente: el montaje se revisa antes de continuar.
\end{tcolorbox}
\end{document}
